\begin{proposition}[径向依赖残差的 Stokes 曲率阻碍与常数反项的不可消去性]\label{prop:xi-radial-stokes-curvature-counterterm-obstruction}
令 \(n\ge 1\)，取圆柱流形 \(\TT^{n}\times I\)（\(I\subset\RR\) 为区间，径向参数记为 \(\rho\)）。
给定可微函数
\[
m:I\to\RR^{n},\qquad \rho\longmapsto (m_1(\rho),\dots,m_n(\rho)),
\]
并定义 \(1\)-形式
\[
A(\rho):=\sum_{j=1}^{n} m_j(\rho)\,\frac{d\phi_j}{2\pi}\ \in\ \Omega^{1}(\TT^{n}\times I),
\]
其中 \(\phi_j\) 为第 \(j\) 个圆周角坐标。
则其曲率 \(2\)-形式为
\[
F:=dA=\sum_{j=1}^{n} m'_j(\rho)\,d\rho\wedge \frac{d\phi_j}{2\pi}.
\]
考虑常数反项（与 \(\mathrm{CT}_\eta\) 的“常数级平移”机制同型，参见备注 \ref{rem:xi-finite-part-drift-torsor}）：
\[
m(\rho)\longmapsto m(\rho)+\eta,\qquad \eta\in\RR^{n}\ \text{常向量}.
\]
则：
\begin{enumerate}
  \item 该常数反项不改变 \(F\)。特别地，若存在某 \(j\) 使 \(m'_j(\rho)\not\equiv 0\)，则任何常数反项都无法使曲率消失。
  \item 对任意矩形 \(R:=\TT^{1}_{j}\times[\rho_0,\rho_1]\subset \TT^{n}\times I\)（\(\TT^{1}_{j}\) 为第 \(j\) 个圆周因子）成立 Stokes 恒等式
  \[
  \boxed{\;
  \int_{R}F=\int_{\partial R}A=m_j(\rho_1)-m_j(\rho_0)\; } .
  \]
  因此径向依赖在上同调意义上等价于 \(F\) 在
  \(H^{2}(\TT^{n}\times I,\TT^{n}\times\partial I;\RR)\)
  中的非平凡性；它刻画了“仅靠常数反项无法抵消径向残差”的结构性障碍。
\end{enumerate}
\end{proposition}
\begin{proof}
常数反项对应
\(
A\mapsto A+\sum_{j}\eta_j\frac{d\phi_j}{2\pi}
\)，其外微分为零，故 \(F=dA\) 不变，得 (1)。
\par
对 (2)，在 \(R\) 上有
\(
F=d\bigl(m_j(\rho)\frac{d\phi_j}{2\pi}\bigr)
\)；
由带角流形上的 Stokes 定理（命题 \ref{prop:cdim-stokes-half-face-count} 的局部模型）得
\[
\int_{R}F=\int_{\partial R}A
=\int_{\TT^1_j}\left(m_j(\rho_1)\frac{d\phi_j}{2\pi}-m_j(\rho_0)\frac{d\phi_j}{2\pi}\right)
=m_j(\rho_1)-m_j(\rho_0).
\]
最后的同调陈述由“曲率为零 \(\Leftrightarrow\) 连接为闭且可被常数反项规约”为平凡类”的标准等价刻画推出。证毕。
\leanverified{paper\_xi\_radial\_stokes\_curvature\_counterterm\_obstruction}
\end{proof}

\endinput
